\documentclass[12pt]{report}
\usepackage[utf8]{inputenc}
\usepackage[spanish]{babel}
\usepackage{geometry}
\usepackage{graphicx}
\usepackage{amsmath}
\usepackage{booktabs}
\usepackage{array}
\usepackage{multirow}
\usepackage{siunitx}
\usepackage{float}
\usepackage{fancyhdr}
\usepackage{hyperref}
\usepackage{listings}
\usepackage{forest}
\usepackage{pdflscape}
\usepackage{adjustbox}
\newcommand{\ligadura}[1]{{\tiny\color{blue!60!black}$\{$#1$\}$}}
\forestset{
  arbolsld/.style={
    for tree={
      font=\scriptsize, align=center, inner sep=1pt,
      l sep=10pt, s sep=6pt, parent anchor=south, child anchor=north,
      edge={thin}
    }
  },
  exito/.style={font=\scriptsize\bfseries, text=green!50!black},
  fallo/.style={font=\scriptsize\bfseries, text=red!70!black},
  corte/.style={},
  negacion/.style={draw, rounded corners=1pt, inner sep=2pt},
}


\geometry{left=2.5cm, right=2.5cm, top=3cm, bottom=3cm}
\graphicspath{{plantilla/}}

\lstset{
    language=Prolog,
    basicstyle=\ttfamily\footnotesize,
    breaklines=true,
    columns=fullflexible,
    keepspaces=true,
    showstringspaces=false,
    frame=single,
    framesep=4pt,
    tabsize=4,
    aboveskip=8pt,
    belowskip=8pt
}

\def \unidad{\textbf{Facultad de Ingenierías}}
\def \programa{Ingeniería de Sistemas y Computación}
\def \curso{IS453}
\def \titulo{ProgIIIG1}
\def \subtitulo {Actividad03 : Casos con Listas}
\def \autores{
    Andrés David Rincón Salazar\\
    andres.rincon2@utp.edu.co\\
    ID: 11.289.05.513\\
}
\def \fecha{- Octubre 2026 -}
\def \lugar{\textbf{Pereira, Colombia}}


\begin{document}
\begin{titlepage}
    \begin{center}
        \vspace*{0.75cm}
        
        \includegraphics[width=0.5\linewidth]{Imágenes/Institucional/Escudo-utp1.jpg}

        \vspace*{0.75cm}
        \LARGE
        \unidad\\
        \programa\\
        \curso
        
        \vspace{0.75cm}
        
        \Huge
        \textbf{\titulo}
            
        \vspace{0.75cm}
        \LARGE
        \subtitulo
            
        \vspace{0.75cm}
        
        \large    
        \autores
            
        %\vfill
        
        \vspace*{0.75cm}
        \lugar\\
        \fecha
        
    \end{center}
\end{titlepage}
\tableofcontents
\newpage


%===============================================================================
% INTRODUCCION
%===============================================================================
\chapter{Introducción}

Esta actividad retoma dos ejercicios ya resueltos con hechos simples y los
reescribe usando \textbf{estructuras} (términos compuestos) y
\textbf{listas}. Las reglas se modificaron para que, aun con la nueva
representación de los datos, el programa resuelva exactamente las mismas
consultas:

\begin{itemize}
    \item \textbf{Punto 1:} familia Simpson (numeral 1 de la Actividad 1).
          Los hechos \texttt{padre\_de/2}, \texttt{madre\_de/2},
          \texttt{hombre/1} y \texttt{mujer/1} se reemplazan por las
          estructuras \texttt{persona(Nombre, Sexo)} y
          \texttt{familia(Progenitores, Hijos)}, donde los dos argumentos de
          \texttt{familia/2} son listas.
    \item \textbf{Punto 2:} ciudades de Canadá (punto 3 de la Actividad 2).
          Los hechos \texttt{arista(Origen, Destino, Costo)} se reemplazan
          por una lista de adyacencia por ciudad,
          \texttt{nodo(Ciudad, [enlace(Destino, Costo), \ldots])}.
\end{itemize}

Cada programa conserva su encabezado descriptivo (propósito, predicados y
salvedades), comentarios de sección y un bloque de consultas de validación.

\section*{Verificación de equivalencia}

Que ``siga resolviendo las mismas consultas'' se comprobó de forma
automática y no solo con las consultas del encabezado. Se cargaron el
programa original y el nuevo en SWI-Prolog 10.0.2 y se compararon sus
respuestas:

\begin{itemize}
    \item \textbf{Simpson:} para cada uno de los 11 predicados de parentesco
          se probaron los $13 \times 13$ pares de individuos, y se
          comparó además la enumeración completa con ambos argumentos sin
          instanciar y el sexo de cada individuo. Los conjuntos de
          respuestas son idénticos. El orden de las soluciones sí puede
          cambiar, porque ahora se recorren listas y no hechos sueltos.
    \item \textbf{Canadá:} se compararon, con el orden de las soluciones,
          \texttt{conexion\_directa/3}, \texttt{conectado\_con/4},
          \texttt{tiene\_aristas/1}, \texttt{camino/4}, \texttt{hay\_camino/2}
          y \texttt{costo\_via/4} sobre todos los pares y ternas de las seis
          ciudades más Toronto. Las respuestas son idénticas, también en
          orden, porque las listas de adyacencia conservan el orden de los
          hechos originales.
\end{itemize}

\section*{Árboles de decisión (SLD)}

Como en la Actividad 2, al final de cada punto se incluye el árbol SLD de
su consulta principal, generado con el paquete \emph{SLDNF Draw} de
SWI-Prolog (\texttt{pack\_install(sldnfdraw)}) a partir de las mismas
cláusulas del programa y dibujado con el paquete \texttt{forest} de
\LaTeX{}. Los archivos de entrada y el script que los regenera están en
\texttt{arboles/fuentes}. Para que la herramienta pueda dibujarlos, en esos
archivos \verb|\+| se escribe \texttt{not/1} (es la misma negación por
fallo), las variables anónimas \texttt{\_} llevan nombre y
\texttt{member/2} se escribe con sus dos cláusulas estándar. Esto último es
deliberado: con listas, el recorrido de \texttt{member/2} es justo lo que
hay que ver en el árbol. La notación es la de la Actividad 2:

\begin{itemize}
    \item Cada nodo es el resolvente (los objetivos pendientes); las
          variables llevan un subíndice que indica en qué paso de
          resolución fueron renombradas.
    \item El texto azul entre llaves es la ligadura (sustitución) que
          produjo la unificación con la cláusula elegida.
    \item Un recuadro es una negación por fallo: dentro está el subárbol
          de \texttt{memberchk/2} y su resultado.
    \item \textbf{true} (verde) es una rama de éxito, \textbf{false} (rojo)
          una rama fallida y $\cdots$ un subárbol que existe pero no se
          dibuja.
    \item Prolog recorre el árbol en profundidad y de izquierda a derecha;
          el \emph{backtracking} es el paso de una rama a la siguiente rama
          hermana a la derecha.
\end{itemize}

Una consecuencia de usar listas: el mismo razonamiento necesita más niveles
que con hechos sueltos, porque cada acceso a un hecho de la Actividad 1 o 2
pasa ahora por un hecho \texttt{familia/2} o \texttt{nodo/2} y por un
recorrido de \texttt{member/2}. Los árboles son anchos o altos; el PDF es
vectorial, así que se pueden ampliar sin perder nitidez.

\newpage


%===============================================================================
% PUNTO 1
%===============================================================================
\chapter{Punto 1: Familia Simpson}

\section{Cambio de representación}

\begin{table}[H]
    \centering
    \small
    \renewcommand{\arraystretch}{1.3}
    \begin{tabular}{>{\raggedright\arraybackslash}p{3.8cm} >{\raggedright\arraybackslash}p{5.3cm} >{\raggedright\arraybackslash}p{5.3cm}}
        \toprule
        \textbf{Concepto} & \textbf{Actividad 1 (hechos)} & \textbf{Actividad 3 (estructuras y listas)} \\
        \midrule
        Padres e hijos & \texttt{padre\_de(homero, bart).}\newline \texttt{madre\_de(marge, bart).}
                       & \texttt{familia([homero, marge], [bart, lisa, maggie]).} \\
        Sexo & \texttt{hombre(herbert).}\newline y reglas \texttt{hombre(N) :- padre\_de(N, \_).}
             & \texttt{persona(herbert, masculino).} \\
        \texttt{padre\_de/2}, \texttt{madre\_de/2} & hechos
             & reglas: \texttt{progenitor\_de/2} filtrado por \texttt{hombre/1} o \texttt{mujer/1} \\
        \texttt{progenitor\_de/2} & dos cláusulas (padre o madre)
             & \texttt{member/2} sobre las listas de una \texttt{familia/2} \\
        Hermanos, abuelos, primos & recorren \texttt{progenitor\_de/2} de a un hijo
             & usan \texttt{hijos\_de/2}, que devuelve la lista de hijos con \texttt{findall/3} \\
        \bottomrule
    \end{tabular}
\end{table}

Decisiones de diseño:

\begin{itemize}
    \item El sexo ya no se deduce de ``ser padre'' o ``ser madre'': se
          declara una vez por persona. Por eso \texttt{padre\_de/2} y
          \texttt{madre\_de/2} pasan a ser reglas, y los hechos
          \texttt{hombre/1} y \texttt{mujer/1} dejan de estar repartidos
          entre hechos y reglas (lo que causaba la advertencia de cláusulas
          no contiguas de la Actividad 1).
    \item Herbert y Ling tienen un solo progenitor conocido, así que cada
          uno va en su propio \texttt{familia/2} con una lista de un
          elemento. Meter a Herbert en la familia de Homero haría a Mona su
          madre, algo que la Actividad 1 nunca afirmó.
    \item Una persona sin hecho \texttt{persona/2} no tiene sexo, y las reglas
          que lo exigen fallan para ella. Cada individuo nuevo debe
          registrarse.
\end{itemize}

\section{Código}

\lstinputlisting[caption={ProgIIIG1-Act03-Pto01-Andres\_David\_Rincon\_Salazar.pl}]{../ProgIIIG1-Act03-Pto01-Andres_David_Rincon_Salazar.pl}

\section{Árbol SLD}

Consulta dibujada: \texttt{padre\_de(P, bart)}. Es la consulta que mejor
muestra el cambio de representación, porque obliga a recorrer las listas de
\texttt{familia/2} y a filtrar por sexo con \texttt{persona/2}. El árbol se
muestra completo (8 niveles) y es finito:

\begin{itemize}
    \item El primer nivel desarrolla \texttt{padre\_de/2} en
          \texttt{progenitor\_de(P, bart), hombre(P)}, y el segundo
          escoge un hecho \texttt{familia/2}. Hay cinco hechos, así que
          hay cinco ramas hermanas.
    \item En cuatro de ellas \texttt{member(bart, Hijos)} recorre la lista
          de hijos hasta llegar a \texttt{[]} sin encontrar a Bart, y la
          rama termina en \texttt{false}. Las familias de Herbert y de Ling
          tienen un solo hijo, por eso esas ramas son cortas.
    \item En \texttt{familia([homero, marge], [bart, lisa, maggie])} Bart
          es la cabeza de la lista, así que la primera cláusula de
          \texttt{member/2} tiene éxito. Luego
          \texttt{member(P, [homero, marge])} produce \texttt{P = homero}
          y \texttt{hombre(homero)} se reduce a
          \texttt{persona(homero, masculino)}: es la única rama de éxito.
    \item Al retroceder, \texttt{P = marge} llega a
          \texttt{persona(marge, masculino)}, que no existe, y falla. Esa
          es la rama donde el filtro \texttt{hombre/1} convierte a
          \texttt{progenitor\_de/2} en \texttt{padre\_de/2}. Después Prolog
          sigue recorriendo \texttt{[lisa, maggie]} en busca de otro
          \texttt{bart}, que no hay.
\end{itemize}

\begin{landscape}
\begin{figure}[p]
    \centering
    \adjustbox{width=\linewidth, totalheight=0.85\textheight, keepaspectratio}{\begin{forest} arbolsld
[{padre\_de($P_{0}$,bart)}
  [{progenitor\_de($P_{0}$,bart),\\hombre($P_{0}$)}
    [{familia($Progenitores_{0}$,$Hijos_{0}$),\\member(bart,$Hijos_{0}$),\\member($P_{0}$,$Progenitores_{0}$),\\hombre($P_{0}$)}
      [{\ligadura{$Progenitores_{0}$/\lbrack abraham,mona\rbrack , $Hijos_{0}$/\lbrack homero\rbrack}\\member(bart,\lbrack homero\rbrack ),\\member($P_{0}$,\lbrack abraham,mona\rbrack ),\\hombre($P_{0}$)}
        [{member(bart,\lbrack\rbrack ),\\member($P_{0}$,\lbrack abraham,mona\rbrack ),\\hombre($P_{0}$)}
          [{false}, fallo
          ]
        ]
      ]
      [{\ligadura{$Progenitores_{0}$/\lbrack abraham\rbrack , $Hijos_{0}$/\lbrack herbert\rbrack}\\member(bart,\lbrack herbert\rbrack ),\\member($P_{0}$,\lbrack abraham\rbrack ),\\hombre($P_{0}$)}
        [{member(bart,\lbrack\rbrack ),\\member($P_{0}$,\lbrack abraham\rbrack ),\\hombre($P_{0}$)}
          [{false}, fallo
          ]
        ]
      ]
      [{\ligadura{$Progenitores_{0}$/\lbrack homero,marge\rbrack , $Hijos_{0}$/\lbrack bart,lisa,maggie\rbrack}\\member(bart,\lbrack bart,lisa,maggie\rbrack ),\\member($P_{0}$,\lbrack homero,marge\rbrack ),\\hombre($P_{0}$)}
        [{member($P_{0}$,\lbrack homero,marge\rbrack ),\\hombre($P_{0}$)}
          [{\ligadura{$P_{0}$/homero}\\hombre(homero)}
            [{persona(homero,masculino)}
              [{true}, exito
              ]
            ]
          ]
          [{member($P_{0}$,\lbrack marge\rbrack ),\\hombre($P_{0}$)}
            [{\ligadura{$P_{0}$/marge}\\hombre(marge)}
              [{persona(marge,masculino)}
                [{false}, fallo
                ]
              ]
            ]
            [{member($P_{0}$,\lbrack\rbrack ),\\hombre($P_{0}$)}
              [{false}, fallo
              ]
            ]
          ]
        ]
        [{member(bart,\lbrack lisa,maggie\rbrack ),\\member($P_{0}$,\lbrack homero,marge\rbrack ),\\hombre($P_{0}$)}
          [{member(bart,\lbrack maggie\rbrack ),\\member($P_{0}$,\lbrack homero,marge\rbrack ),\\hombre($P_{0}$)}
            [{member(bart,\lbrack\rbrack ),\\member($P_{0}$,\lbrack homero,marge\rbrack ),\\hombre($P_{0}$)}
              [{false}, fallo
              ]
            ]
          ]
        ]
      ]
      [{\ligadura{$Progenitores_{0}$/\lbrack clancy,jacqueline\rbrack , $Hijos_{0}$/\lbrack patty,selma,marge\rbrack}\\member(bart,\lbrack patty,selma,marge\rbrack ),\\member($P_{0}$,\lbrack clancy,jacqueline\rbrack ),\\hombre($P_{0}$)}
        [{member(bart,\lbrack selma,marge\rbrack ),\\member($P_{0}$,\lbrack clancy,jacqueline\rbrack ),\\hombre($P_{0}$)}
          [{member(bart,\lbrack marge\rbrack ),\\member($P_{0}$,\lbrack clancy,jacqueline\rbrack ),\\hombre($P_{0}$)}
            [{member(bart,\lbrack\rbrack ),\\member($P_{0}$,\lbrack clancy,jacqueline\rbrack ),\\hombre($P_{0}$)}
              [{false}, fallo
              ]
            ]
          ]
        ]
      ]
      [{\ligadura{$Progenitores_{0}$/\lbrack selma\rbrack , $Hijos_{0}$/\lbrack ling\rbrack}\\member(bart,\lbrack ling\rbrack ),\\member($P_{0}$,\lbrack selma\rbrack ),\\hombre($P_{0}$)}
        [{member(bart,\lbrack\rbrack ),\\member($P_{0}$,\lbrack selma\rbrack ),\\hombre($P_{0}$)}
          [{false}, fallo
          ]
        ]
      ]
    ]
  ]
]
\end{forest}
}
    \caption{Árbol SLD completo de \texttt{padre\_de(P, bart)}.}
\end{figure}
\end{landscape}

\newpage


%===============================================================================
% PUNTO 2
%===============================================================================
\chapter{Punto 2: Ciudades de Canadá}

\section{Cambio de representación}

\begin{table}[H]
    \centering
    \small
    \renewcommand{\arraystretch}{1.3}
    \begin{tabular}{>{\raggedright\arraybackslash}p{3.8cm} >{\raggedright\arraybackslash}p{5.3cm} >{\raggedright\arraybackslash}p{5.3cm}}
        \toprule
        \textbf{Concepto} & \textbf{Actividad 2 (hechos)} & \textbf{Actividad 3 (estructuras y listas)} \\
        \midrule
        Aristas & nueve hechos \texttt{arista(Origen, Destino, Costo).}
                & seis hechos \texttt{nodo(Ciudad, Enlaces)}; \texttt{Enlaces} es una lista de estructuras \texttt{enlace(Destino, Costo)} \\
        Arista directa & hecho \texttt{arista/3}
                       & \texttt{conexion\_directa/3}: \texttt{nodo(A, Es), member(enlace(B, C), Es)} \\
        \texttt{conectado\_con/4}, \texttt{tiene\_aristas/1}, \texttt{recorrer/5}, \texttt{ruta\_ingenua/2}
                       & usan \texttt{arista/3}
                       & usan \texttt{conexion\_directa/3} \\
        Vecinos de un nodo & no existía
                           & \texttt{vecinos\_de/2} devuelve la lista con \texttt{findall/3} \\
        \bottomrule
    \end{tabular}
\end{table}

Decisiones de diseño:

\begin{itemize}
    \item El predicado \texttt{arista/3} desaparece. Todo acceso a las aristas
          pasa por \texttt{conexion\_directa/3}, que es la única regla que
          recorre las listas; el resto del programa no depende de cómo se
          guardan.
    \item Winnipeg no tiene aristas de salida pero sí debe tener su hecho:
          \texttt{nodo(winnipeg, [])}. Si faltara, \texttt{tiene\_aristas/1}
          seguiría respondiendo \texttt{true} (porque Winnipeg es destino),
          pero \texttt{vecinos\_de(winnipeg, V)} fallaría en lugar de
          responder \texttt{V = []}.
    \item Se conservó \texttt{ruta\_ingenua/2} porque la Actividad 2 la usa
          para mostrar que el árbol SLD es infinito por el ciclo
          Edmonton $\to$ Saskatoon $\to$ Calgary $\to$ Edmonton. La versión
          segura (\texttt{camino/4}, \texttt{hay\_camino/2}) sigue usando
          una lista de nodos visitados.
\end{itemize}

\section{Código}

\lstinputlisting[caption={ProgIIIG1-Act03-Pto02-Andres\_David\_Rincon\_Salazar.pl}]{../ProgIIIG1-Act03-Pto02-Andres_David_Rincon_Salazar.pl}

\section{Árboles SLD}

\subsection{Versión segura (árbol completo)}

\begin{sloppypar}
Consulta dibujada: \texttt{recorrer(edmonton, calgary, [edmonton], R, C)},
que es el núcleo de \texttt{camino(edmonton, calgary, Nodos, Costo)}. El
árbol se muestra completo y es finito:
\end{sloppypar}

\begin{itemize}
    \item Tiene una sola rama de éxito, con
          \texttt{R = [calgary, saskatoon, edmonton]} y
          \texttt{C = 21}. \texttt{camino/4} invierte la lista y
          responde \texttt{[edmonton, saskatoon, calgary]}.
    \item Cada paso por el grafo se desdobla en tres niveles:
          \texttt{conexion\_directa/3}, el hecho \texttt{nodo/2} que
          entrega la lista de enlaces y el recorrido de esa lista con
          \texttt{member/2}. En la Actividad 2 el mismo paso era un solo
          nivel (\texttt{arista/3}). Cuando la lista tiene más de un
          enlace, la segunda cláusula de \texttt{member/2} produce la rama
          hermana con el resto de la lista (por ejemplo
          \texttt{enlace(regina, 14)} después de
          \texttt{enlace(edmonton, 4)}).
    \item Al seguir buscando desde Calgary, el enlace hacia Edmonton se
          poda porque \texttt{memberchk(edmonton, ...)} es verdadero: la
          lista de visitados corta el ciclo que hace infinita a la versión
          ingenua.
    \item Las demás ramas terminan en \texttt{false}: o vuelven a un nodo
          ya visitado, o llegan a Winnipeg, cuya lista de enlaces es
          \texttt{[]}, de modo que \texttt{member/2} falla de inmediato.
\end{itemize}
\begin{figure}[p]
    \centering
    \adjustbox{width=\linewidth, totalheight=0.9\textheight, keepaspectratio}{\begin{forest} arbolsld
[{recorrer(edmonton,calgary,\lbrack edmonton\rbrack ,$R_{0}$,$C_{0}$)}
  [{conexion\_directa(edmonton,$Siguiente_{0}$,$C1_{0}$),\\not(memberchk($Siguiente_{0}$,\lbrack edmonton\rbrack )),\\recorrer($Siguiente_{0}$,calgary,\lbrack $Siguiente_{0}$,edmonton\rbrack ,$R_{0}$,$C2_{0}$),\\$C_{0}$ is $C1_{0}$+$C2_{0}$}
    [{nodo(edmonton,$Enlaces_{0}$),\\member(enlace($Siguiente_{0}$,$C1_{0}$),$Enlaces_{0}$),\\not(memberchk($Siguiente_{0}$,\lbrack edmonton\rbrack )),\\recorrer($Siguiente_{0}$,calgary,\lbrack $Siguiente_{0}$,edmonton\rbrack ,$R_{0}$,$C2_{0}$),\\$C_{0}$ is $C1_{0}$+$C2_{0}$}
      [{\ligadura{$Enlaces_{0}$/\lbrack enlace(saskatoon,12)\rbrack}\\member(enlace($Siguiente_{0}$,$C1_{0}$),\lbrack enlace(saskatoon,12)\rbrack ),\\not(memberchk($Siguiente_{0}$,\lbrack edmonton\rbrack )),\\recorrer($Siguiente_{0}$,calgary,\lbrack $Siguiente_{0}$,edmonton\rbrack ,$R_{0}$,$C2_{0}$),\\$C_{0}$ is $C1_{0}$+$C2_{0}$}
        [{\ligadura{$Siguiente_{0}$/saskatoon, $C1_{0}$/12}\\not(memberchk(saskatoon,\lbrack edmonton\rbrack )),\\recorrer(saskatoon,calgary,\lbrack saskatoon,edmonton\rbrack ,$R_{0}$,$C2_{0}$),\\$C_{0}$ is 12+$C2_{0}$}
          [{memberchk(saskatoon,\lbrack edmonton\rbrack )\\$\Rightarrow$ false}, negacion
            [{recorrer(saskatoon,calgary,\lbrack saskatoon,edmonton\rbrack ,$R_{0}$,$C2_{0}$),\\$C_{0}$ is 12+$C2_{0}$}
              [{conexion\_directa(saskatoon,$Siguiente_{1}$,$C1_{1}$),\\not(memberchk($Siguiente_{1}$,\lbrack saskatoon,edmonton\rbrack )),\\recorrer($Siguiente_{1}$,calgary,\lbrack $Siguiente_{1}$,saskatoon,edmonton\rbrack ,$R_{0}$,$C2_{1}$),\\$C2_{0}$ is $C1_{1}$+$C2_{1}$,\\$C_{0}$ is 12+$C2_{0}$}
                [{nodo(saskatoon,$Enlaces_{1}$),\\member(enlace($Siguiente_{1}$,$C1_{1}$),$Enlaces_{1}$),\\not(memberchk($Siguiente_{1}$,\lbrack saskatoon,edmonton\rbrack )),\\recorrer($Siguiente_{1}$,calgary,\lbrack $Siguiente_{1}$,saskatoon,edmonton\rbrack ,$R_{0}$,$C2_{1}$),\\$C2_{0}$ is $C1_{1}$+$C2_{1}$,\\$C_{0}$ is 12+$C2_{0}$}
                  [{\ligadura{$Enlaces_{1}$/\lbrack enlace(calgary,9),enlace(winnipeg,20)\rbrack}\\member(enlace($Siguiente_{1}$,$C1_{1}$),\lbrack enlace(calgary,9),enlace(winnipeg,20)\rbrack ),\\not(memberchk($Siguiente_{1}$,\lbrack saskatoon,edmonton\rbrack )),\\recorrer($Siguiente_{1}$,calgary,\lbrack $Siguiente_{1}$,saskatoon,edmonton\rbrack ,$R_{0}$,$C2_{1}$),\\$C2_{0}$ is $C1_{1}$+$C2_{1}$,\\$C_{0}$ is 12+$C2_{0}$}
                    [{\ligadura{$Siguiente_{1}$/calgary, $C1_{1}$/9}\\not(memberchk(calgary,\lbrack saskatoon,edmonton\rbrack )),\\recorrer(calgary,calgary,\lbrack calgary,saskatoon,edmonton\rbrack ,$R_{0}$,$C2_{1}$),\\$C2_{0}$ is 9+$C2_{1}$,\\$C_{0}$ is 12+$C2_{0}$}
                      [{memberchk(calgary,\lbrack saskatoon,edmonton\rbrack )\\$\Rightarrow$ false}, negacion
                        [{recorrer(calgary,calgary,\lbrack calgary,saskatoon,edmonton\rbrack ,$R_{0}$,$C2_{1}$),\\$C2_{0}$ is 9+$C2_{1}$,\\$C_{0}$ is 12+$C2_{0}$}
                          [{\ligadura{$R_{0}$/\lbrack calgary,saskatoon,edmonton\rbrack , $C2_{1}$/0}\\$C2_{0}$ is 9+0,\\$C_{0}$ is 12+$C2_{0}$}
                            [{\ligadura{$C2_{0}$/9}\\$C_{0}$ is 12+9}
                              [{\ligadura{$C_{0}$/21}\\true}, exito
                              ]
                            ]
                          ]
                          [{conexion\_directa(calgary,$Siguiente_{2}$,$C1_{2}$),\\not(memberchk($Siguiente_{2}$,\lbrack calgary,saskatoon,edmonton\rbrack )),\\recorrer($Siguiente_{2}$,calgary,\lbrack $Siguiente_{2}$,calgary,saskatoon,edmonton\rbrack ,$R_{0}$,$C2_{2}$),\\$C2_{1}$ is $C1_{2}$+$C2_{2}$,\\$C2_{0}$ is 9+$C2_{1}$,\\$C_{0}$ is 12+$C2_{0}$}
                            [{nodo(calgary,$Enlaces_{2}$),\\member(enlace($Siguiente_{2}$,$C1_{2}$),$Enlaces_{2}$),\\not(memberchk($Siguiente_{2}$,\lbrack calgary,saskatoon,edmonton\rbrack )),\\recorrer($Siguiente_{2}$,calgary,\lbrack $Siguiente_{2}$,calgary,saskatoon,edmonton\rbrack ,$R_{0}$,$C2_{2}$),\\$C2_{1}$ is $C1_{2}$+$C2_{2}$,\\$C2_{0}$ is 9+$C2_{1}$,\\$C_{0}$ is 12+$C2_{0}$}
                              [{\ligadura{$Enlaces_{2}$/\lbrack enlace(edmonton,4),enlace(regina,14)\rbrack}\\member(enlace($Siguiente_{2}$,$C1_{2}$),\lbrack enlace(edmonton,4),enlace(regina,14)\rbrack ),\\not(memberchk($Siguiente_{2}$,\lbrack calgary,saskatoon,edmonton\rbrack )),\\recorrer($Siguiente_{2}$,calgary,\lbrack $Siguiente_{2}$,calgary,saskatoon,edmonton\rbrack ,$R_{0}$,$C2_{2}$),\\$C2_{1}$ is $C1_{2}$+$C2_{2}$,\\$C2_{0}$ is 9+$C2_{1}$,\\$C_{0}$ is 12+$C2_{0}$}
                                [{\ligadura{$Siguiente_{2}$/edmonton, $C1_{2}$/4}\\not(memberchk(edmonton,\lbrack calgary,saskatoon,edmonton\rbrack )),\\recorrer(edmonton,calgary,\lbrack edmonton,calgary,saskatoon,edmonton\rbrack ,$R_{0}$,$C2_{2}$),\\$C2_{1}$ is 4+$C2_{2}$,\\$C2_{0}$ is 9+$C2_{1}$,\\$C_{0}$ is 12+$C2_{0}$}
                                  [{memberchk(edmonton,\lbrack calgary,saskatoon,edmonton\rbrack )\\$\Rightarrow$ true}, negacion
                                    [{false}, fallo
                                    ]
                                  ]
                                ]
                                [{member(enlace($Siguiente_{2}$,$C1_{2}$),\lbrack enlace(regina,14)\rbrack ),\\not(memberchk($Siguiente_{2}$,\lbrack calgary,saskatoon,edmonton\rbrack )),\\recorrer($Siguiente_{2}$,calgary,\lbrack $Siguiente_{2}$,calgary,saskatoon,edmonton\rbrack ,$R_{0}$,$C2_{2}$),\\$C2_{1}$ is $C1_{2}$+$C2_{2}$,\\$C2_{0}$ is 9+$C2_{1}$,\\$C_{0}$ is 12+$C2_{0}$}
                                  [{\ligadura{$Siguiente_{2}$/regina, $C1_{2}$/14}\\not(memberchk(regina,\lbrack calgary,saskatoon,edmonton\rbrack )),\\recorrer(regina,calgary,\lbrack regina,calgary,saskatoon,edmonton\rbrack ,$R_{0}$,$C2_{2}$),\\$C2_{1}$ is 14+$C2_{2}$,\\$C2_{0}$ is 9+$C2_{1}$,\\$C_{0}$ is 12+$C2_{0}$}
                                    [{memberchk(regina,\lbrack calgary,saskatoon,edmonton\rbrack )\\$\Rightarrow$ false}, negacion
                                      [{recorrer(regina,calgary,\lbrack regina,calgary,saskatoon,edmonton\rbrack ,$R_{0}$,$C2_{2}$),\\$C2_{1}$ is 14+$C2_{2}$,\\$C2_{0}$ is 9+$C2_{1}$,\\$C_{0}$ is 12+$C2_{0}$}
                                        [{conexion\_directa(regina,$Siguiente_{3}$,$C1_{3}$),\\not(memberchk($Siguiente_{3}$,\lbrack regina,calgary,saskatoon,edmonton\rbrack )),\\recorrer($Siguiente_{3}$,calgary,\lbrack $Siguiente_{3}$,regina,calgary,saskatoon,edmonton\rbrack ,$R_{0}$,$C2_{3}$),\\$C2_{2}$ is $C1_{3}$+$C2_{3}$,\\$C2_{1}$ is 14+$C2_{2}$,\\$C2_{0}$ is 9+$C2_{1}$,\\$C_{0}$ is 12+$C2_{0}$}
                                          [{nodo(regina,$Enlaces_{3}$),\\member(enlace($Siguiente_{3}$,$C1_{3}$),$Enlaces_{3}$),\\not(memberchk($Siguiente_{3}$,\lbrack regina,calgary,saskatoon,edmonton\rbrack )),\\recorrer($Siguiente_{3}$,calgary,\lbrack $Siguiente_{3}$,regina,calgary,saskatoon,edmonton\rbrack ,$R_{0}$,$C2_{3}$),\\$C2_{2}$ is $C1_{3}$+$C2_{3}$,\\$C2_{1}$ is 14+$C2_{2}$,\\$C2_{0}$ is 9+$C2_{1}$,\\$C_{0}$ is 12+$C2_{0}$}
                                            [{\ligadura{$Enlaces_{3}$/\lbrack enlace(saskatoon,7),enlace(winnipeg,4)\rbrack}\\member(enlace($Siguiente_{3}$,$C1_{3}$),\lbrack enlace(saskatoon,7),enlace(winnipeg,4)\rbrack ),\\not(memberchk($Siguiente_{3}$,\lbrack regina,calgary,saskatoon,edmonton\rbrack )),\\recorrer($Siguiente_{3}$,calgary,\lbrack $Siguiente_{3}$,regina,calgary,saskatoon,edmonton\rbrack ,$R_{0}$,$C2_{3}$),\\$C2_{2}$ is $C1_{3}$+$C2_{3}$,\\$C2_{1}$ is 14+$C2_{2}$,\\$C2_{0}$ is 9+$C2_{1}$,\\$C_{0}$ is 12+$C2_{0}$}
                                              [{\ligadura{$Siguiente_{3}$/saskatoon, $C1_{3}$/7}\\not(memberchk(saskatoon,\lbrack regina,calgary,saskatoon,edmonton\rbrack )),\\recorrer(saskatoon,calgary,\lbrack saskatoon,regina,calgary,saskatoon,edmonton\rbrack ,$R_{0}$,$C2_{3}$),\\$C2_{2}$ is 7+$C2_{3}$,\\$C2_{1}$ is 14+$C2_{2}$,\\$C2_{0}$ is 9+$C2_{1}$,\\$C_{0}$ is 12+$C2_{0}$}
                                                [{memberchk(saskatoon,\lbrack regina,calgary,saskatoon,edmonton\rbrack )\\$\Rightarrow$ true}, negacion
                                                  [{false}, fallo
                                                  ]
                                                ]
                                              ]
                                              [{member(enlace($Siguiente_{3}$,$C1_{3}$),\lbrack enlace(winnipeg,4)\rbrack ),\\not(memberchk($Siguiente_{3}$,\lbrack regina,calgary,saskatoon,edmonton\rbrack )),\\recorrer($Siguiente_{3}$,calgary,\lbrack $Siguiente_{3}$,regina,calgary,saskatoon,edmonton\rbrack ,$R_{0}$,$C2_{3}$),\\$C2_{2}$ is $C1_{3}$+$C2_{3}$,\\$C2_{1}$ is 14+$C2_{2}$,\\$C2_{0}$ is 9+$C2_{1}$,\\$C_{0}$ is 12+$C2_{0}$}
                                                [{\ligadura{$Siguiente_{3}$/winnipeg, $C1_{3}$/4}\\not(memberchk(winnipeg,\lbrack regina,calgary,saskatoon,edmonton\rbrack )),\\recorrer(winnipeg,calgary,\lbrack winnipeg,regina,calgary,saskatoon,edmonton\rbrack ,$R_{0}$,$C2_{3}$),\\$C2_{2}$ is 4+$C2_{3}$,\\$C2_{1}$ is 14+$C2_{2}$,\\$C2_{0}$ is 9+$C2_{1}$,\\$C_{0}$ is 12+$C2_{0}$}
                                                  [{memberchk(winnipeg,\lbrack regina,calgary,saskatoon,edmonton\rbrack )\\$\Rightarrow$ false}, negacion
                                                    [{recorrer(winnipeg,calgary,\lbrack winnipeg,regina,calgary,saskatoon,edmonton\rbrack ,$R_{0}$,$C2_{3}$),\\$C2_{2}$ is 4+$C2_{3}$,\\$C2_{1}$ is 14+$C2_{2}$,\\$C2_{0}$ is 9+$C2_{1}$,\\$C_{0}$ is 12+$C2_{0}$}
                                                      [{conexion\_directa(winnipeg,$Siguiente_{4}$,$C1_{4}$),\\not(memberchk($Siguiente_{4}$,\lbrack winnipeg,regina,calgary,saskatoon,edmonton\rbrack )),\\recorrer($Siguiente_{4}$,calgary,\lbrack $Siguiente_{4}$,winnipeg,regina,calgary,saskatoon,edmonton\rbrack ,$R_{0}$,$C2_{4}$),\\$C2_{3}$ is $C1_{4}$+$C2_{4}$,\\$C2_{2}$ is 4+$C2_{3}$,\\$C2_{1}$ is 14+$C2_{2}$,\\$C2_{0}$ is 9+$C2_{1}$,\\$C_{0}$ is 12+$C2_{0}$}
                                                        [{nodo(winnipeg,$Enlaces_{4}$),\\member(enlace($Siguiente_{4}$,$C1_{4}$),$Enlaces_{4}$),\\not(memberchk($Siguiente_{4}$,\lbrack winnipeg,regina,calgary,saskatoon,edmonton\rbrack )),\\recorrer($Siguiente_{4}$,calgary,\lbrack $Siguiente_{4}$,winnipeg,regina,calgary,saskatoon,edmonton\rbrack ,$R_{0}$,$C2_{4}$),\\$C2_{3}$ is $C1_{4}$+$C2_{4}$,\\$C2_{2}$ is 4+$C2_{3}$,\\$C2_{1}$ is 14+$C2_{2}$,\\$C2_{0}$ is 9+$C2_{1}$,\\$C_{0}$ is 12+$C2_{0}$}
                                                          [{\ligadura{$Enlaces_{4}$/\lbrack\rbrack}\\member(enlace($Siguiente_{4}$,$C1_{4}$),\lbrack\rbrack ),\\not(memberchk($Siguiente_{4}$,\lbrack winnipeg,regina,calgary,saskatoon,edmonton\rbrack )),\\recorrer($Siguiente_{4}$,calgary,\lbrack $Siguiente_{4}$,winnipeg,regina,calgary,saskatoon,edmonton\rbrack ,$R_{0}$,$C2_{4}$),\\$C2_{3}$ is $C1_{4}$+$C2_{4}$,\\$C2_{2}$ is 4+$C2_{3}$,\\$C2_{1}$ is 14+$C2_{2}$,\\$C2_{0}$ is 9+$C2_{1}$,\\$C_{0}$ is 12+$C2_{0}$}
                                                            [{false}, fallo
                                                            ]
                                                          ]
                                                        ]
                                                      ]
                                                    ]
                                                  ]
                                                ]
                                                [{member(enlace($Siguiente_{3}$,$C1_{3}$),\lbrack\rbrack ),\\not(memberchk($Siguiente_{3}$,\lbrack regina,calgary,saskatoon,edmonton\rbrack )),\\recorrer($Siguiente_{3}$,calgary,\lbrack $Siguiente_{3}$,regina,calgary,saskatoon,edmonton\rbrack ,$R_{0}$,$C2_{3}$),\\$C2_{2}$ is $C1_{3}$+$C2_{3}$,\\$C2_{1}$ is 14+$C2_{2}$,\\$C2_{0}$ is 9+$C2_{1}$,\\$C_{0}$ is 12+$C2_{0}$}
                                                  [{false}, fallo
                                                  ]
                                                ]
                                              ]
                                            ]
                                          ]
                                        ]
                                      ]
                                    ]
                                  ]
                                  [{member(enlace($Siguiente_{2}$,$C1_{2}$),\lbrack\rbrack ),\\not(memberchk($Siguiente_{2}$,\lbrack calgary,saskatoon,edmonton\rbrack )),\\recorrer($Siguiente_{2}$,calgary,\lbrack $Siguiente_{2}$,calgary,saskatoon,edmonton\rbrack ,$R_{0}$,$C2_{2}$),\\$C2_{1}$ is $C1_{2}$+$C2_{2}$,\\$C2_{0}$ is 9+$C2_{1}$,\\$C_{0}$ is 12+$C2_{0}$}
                                    [{false}, fallo
                                    ]
                                  ]
                                ]
                              ]
                            ]
                          ]
                        ]
                      ]
                    ]
                    [{member(enlace($Siguiente_{1}$,$C1_{1}$),\lbrack enlace(winnipeg,20)\rbrack ),\\not(memberchk($Siguiente_{1}$,\lbrack saskatoon,edmonton\rbrack )),\\recorrer($Siguiente_{1}$,calgary,\lbrack $Siguiente_{1}$,saskatoon,edmonton\rbrack ,$R_{0}$,$C2_{1}$),\\$C2_{0}$ is $C1_{1}$+$C2_{1}$,\\$C_{0}$ is 12+$C2_{0}$}
                      [{\ligadura{$Siguiente_{1}$/winnipeg, $C1_{1}$/20}\\not(memberchk(winnipeg,\lbrack saskatoon,edmonton\rbrack )),\\recorrer(winnipeg,calgary,\lbrack winnipeg,saskatoon,edmonton\rbrack ,$R_{0}$,$C2_{1}$),\\$C2_{0}$ is 20+$C2_{1}$,\\$C_{0}$ is 12+$C2_{0}$}
                        [{memberchk(winnipeg,\lbrack saskatoon,edmonton\rbrack )\\$\Rightarrow$ false}, negacion
                          [{recorrer(winnipeg,calgary,\lbrack winnipeg,saskatoon,edmonton\rbrack ,$R_{0}$,$C2_{1}$),\\$C2_{0}$ is 20+$C2_{1}$,\\$C_{0}$ is 12+$C2_{0}$}
                            [{conexion\_directa(winnipeg,$Siguiente_{2}$,$C1_{2}$),\\not(memberchk($Siguiente_{2}$,\lbrack winnipeg,saskatoon,edmonton\rbrack )),\\recorrer($Siguiente_{2}$,calgary,\lbrack $Siguiente_{2}$,winnipeg,saskatoon,edmonton\rbrack ,$R_{0}$,$C2_{2}$),\\$C2_{1}$ is $C1_{2}$+$C2_{2}$,\\$C2_{0}$ is 20+$C2_{1}$,\\$C_{0}$ is 12+$C2_{0}$}
                              [{nodo(winnipeg,$Enlaces_{2}$),\\member(enlace($Siguiente_{2}$,$C1_{2}$),$Enlaces_{2}$),\\not(memberchk($Siguiente_{2}$,\lbrack winnipeg,saskatoon,edmonton\rbrack )),\\recorrer($Siguiente_{2}$,calgary,\lbrack $Siguiente_{2}$,winnipeg,saskatoon,edmonton\rbrack ,$R_{0}$,$C2_{2}$),\\$C2_{1}$ is $C1_{2}$+$C2_{2}$,\\$C2_{0}$ is 20+$C2_{1}$,\\$C_{0}$ is 12+$C2_{0}$}
                                [{\ligadura{$Enlaces_{2}$/\lbrack\rbrack}\\member(enlace($Siguiente_{2}$,$C1_{2}$),\lbrack\rbrack ),\\not(memberchk($Siguiente_{2}$,\lbrack winnipeg,saskatoon,edmonton\rbrack )),\\recorrer($Siguiente_{2}$,calgary,\lbrack $Siguiente_{2}$,winnipeg,saskatoon,edmonton\rbrack ,$R_{0}$,$C2_{2}$),\\$C2_{1}$ is $C1_{2}$+$C2_{2}$,\\$C2_{0}$ is 20+$C2_{1}$,\\$C_{0}$ is 12+$C2_{0}$}
                                  [{false}, fallo
                                  ]
                                ]
                              ]
                            ]
                          ]
                        ]
                      ]
                      [{member(enlace($Siguiente_{1}$,$C1_{1}$),\lbrack\rbrack ),\\not(memberchk($Siguiente_{1}$,\lbrack saskatoon,edmonton\rbrack )),\\recorrer($Siguiente_{1}$,calgary,\lbrack $Siguiente_{1}$,saskatoon,edmonton\rbrack ,$R_{0}$,$C2_{1}$),\\$C2_{0}$ is $C1_{1}$+$C2_{1}$,\\$C_{0}$ is 12+$C2_{0}$}
                        [{false}, fallo
                        ]
                      ]
                    ]
                  ]
                ]
              ]
            ]
          ]
        ]
        [{member(enlace($Siguiente_{0}$,$C1_{0}$),\lbrack\rbrack ),\\not(memberchk($Siguiente_{0}$,\lbrack edmonton\rbrack )),\\recorrer($Siguiente_{0}$,calgary,\lbrack $Siguiente_{0}$,edmonton\rbrack ,$R_{0}$,$C2_{0}$),\\$C_{0}$ is $C1_{0}$+$C2_{0}$}
          [{false}, fallo
          ]
        ]
      ]
    ]
  ]
]
\end{forest}
}
    \caption{Árbol SLD completo de \texttt{recorrer(edmonton, calgary, [edmonton], R, C)}.}
\end{figure}

\subsection{Versión ingenua (árbol infinito)}

Consulta dibujada: \texttt{ruta\_ingenua(edmonton, calgary)}. Este árbol
es infinito, así que se corta en el nivel 13, donde el objetivo inicial
\texttt{ruta\_ingenua(edmonton, calgary)} reaparece en la rama izquierda
tras una vuelta completa al ciclo \texttt{edmonton $\to$ saskatoon
$\to$ calgary $\to$ edmonton}:

\begin{itemize}
    \item La primera cláusula de \texttt{ruta\_ingenua/2} falla de
          inmediato (Edmonton no tiene enlace directo a Calgary: su lista
          es \texttt{[enlace(saskatoon, 12)]}). La segunda avanza a
          Saskatoon.
    \item En el nivel 8 aparece la primera rama de éxito:
          \texttt{ruta\_ingenua(saskatoon, calgary)} se resuelve con
          \texttt{enlace(calgary, 9)}, la cabeza de la lista de Saskatoon.
          Es el primer \texttt{true} que imprime Prolog.
    \item En el nivel 12 el objetivo inicial vuelve a aparecer. Cada vuelta
          al ciclo cuesta ahora 12 niveles y no 3 como en la Actividad 2,
          porque cada arista se desarrolla en \texttt{conexion\_directa/3},
          \texttt{nodo/2} y \texttt{member/2}; el comportamiento es el
          mismo: un \texttt{true} por vuelta, sin terminar nunca.
    \item Las ramas de la derecha (por Regina y Winnipeg) existen en el
          árbol, pero Prolog nunca llega a ellas: la búsqueda en
          profundidad se queda para siempre en la rama infinita de la
          izquierda. Por eso la estrategia de búsqueda de Prolog (en
          profundidad) no es completa.
\end{itemize}
\begin{landscape}
\begin{figure}[p]
    \centering
    \adjustbox{width=\linewidth, totalheight=0.85\textheight, keepaspectratio}{\begin{forest} arbolsld
[{ruta\_ingenua(edmonton,calgary)}
  [{conexion\_directa(edmonton,calgary,$C_{0}$)}
    [{nodo(edmonton,$Enlaces_{0}$),\\member(enlace(calgary,$C_{0}$),$Enlaces_{0}$)}
      [{\ligadura{$Enlaces_{0}$/\lbrack enlace(saskatoon,12)\rbrack}\\member(enlace(calgary,$C_{0}$),\lbrack enlace(saskatoon,12)\rbrack )}
        [{member(enlace(calgary,$C_{0}$),\lbrack\rbrack )}
          [{false}, fallo
          ]
        ]
      ]
    ]
  ]
  [{conexion\_directa(edmonton,$Intermedio_{0}$,$C_{0}$),\\ruta\_ingenua($Intermedio_{0}$,calgary)}
    [{nodo(edmonton,$Enlaces_{0}$),\\member(enlace($Intermedio_{0}$,$C_{0}$),$Enlaces_{0}$),\\ruta\_ingenua($Intermedio_{0}$,calgary)}
      [{\ligadura{$Enlaces_{0}$/\lbrack enlace(saskatoon,12)\rbrack}\\member(enlace($Intermedio_{0}$,$C_{0}$),\lbrack enlace(saskatoon,12)\rbrack ),\\ruta\_ingenua($Intermedio_{0}$,calgary)}
        [{\ligadura{$Intermedio_{0}$/saskatoon, $C_{0}$/12}\\ruta\_ingenua(saskatoon,calgary)}
          [{conexion\_directa(saskatoon,calgary,$C_{1}$)}
            [{nodo(saskatoon,$Enlaces_{1}$),\\member(enlace(calgary,$C_{1}$),$Enlaces_{1}$)}
              [{\ligadura{$Enlaces_{1}$/\lbrack enlace(calgary,9),enlace(winnipeg,20)\rbrack}\\member(enlace(calgary,$C_{1}$),\lbrack enlace(calgary,9),enlace(winnipeg,20)\rbrack )}
                [{\ligadura{$C_{1}$/9}\\true}, exito
                ]
                [{member(enlace(calgary,$C_{1}$),\lbrack enlace(winnipeg,20)\rbrack )}
                  [{member(enlace(calgary,$C_{1}$),\lbrack\rbrack )}
                    [{false}, fallo
                    ]
                  ]
                ]
              ]
            ]
          ]
          [{conexion\_directa(saskatoon,$Intermedio_{1}$,$C_{1}$),\\ruta\_ingenua($Intermedio_{1}$,calgary)}
            [{nodo(saskatoon,$Enlaces_{1}$),\\member(enlace($Intermedio_{1}$,$C_{1}$),$Enlaces_{1}$),\\ruta\_ingenua($Intermedio_{1}$,calgary)}
              [{\ligadura{$Enlaces_{1}$/\lbrack enlace(calgary,9),enlace(winnipeg,20)\rbrack}\\member(enlace($Intermedio_{1}$,$C_{1}$),\lbrack enlace(calgary,9),enlace(winnipeg,20)\rbrack ),\\ruta\_ingenua($Intermedio_{1}$,calgary)}
                [{\ligadura{$Intermedio_{1}$/calgary, $C_{1}$/9}\\ruta\_ingenua(calgary,calgary)}
                  [{conexion\_directa(calgary,calgary,$C_{2}$)}
                    [{nodo(calgary,$Enlaces_{2}$),\\member(enlace(calgary,$C_{2}$),$Enlaces_{2}$)}
                      [{\ligadura{$Enlaces_{2}$/\lbrack enlace(edmonton,4),enlace(regina,14)\rbrack}\\member(enlace(calgary,$C_{2}$),\lbrack enlace(edmonton,4),enlace(regina,14)\rbrack )}
                        [{member(enlace(calgary,$C_{2}$),\lbrack enlace(regina,14)\rbrack )}
                          [{$\cdots$}, corte]
                        ]
                      ]
                    ]
                  ]
                  [{conexion\_directa(calgary,$Intermedio_{2}$,$C_{2}$),\\ruta\_ingenua($Intermedio_{2}$,calgary)}
                    [{nodo(calgary,$Enlaces_{2}$),\\member(enlace($Intermedio_{2}$,$C_{2}$),$Enlaces_{2}$),\\ruta\_ingenua($Intermedio_{2}$,calgary)}
                      [{\ligadura{$Enlaces_{2}$/\lbrack enlace(edmonton,4),enlace(regina,14)\rbrack}\\member(enlace($Intermedio_{2}$,$C_{2}$),\lbrack enlace(edmonton,4),enlace(regina,14)\rbrack ),\\ruta\_ingenua($Intermedio_{2}$,calgary)}
                        [{\ligadura{$Intermedio_{2}$/edmonton, $C_{2}$/4}\\ruta\_ingenua(edmonton,calgary)}
                          [{$\cdots$}, corte]
                          [{$\cdots$}, corte]
                        ]
                        [{member(enlace($Intermedio_{2}$,$C_{2}$),\lbrack enlace(regina,14)\rbrack ),\\ruta\_ingenua($Intermedio_{2}$,calgary)}
                          [{$\cdots$}, corte]
                          [{$\cdots$}, corte]
                        ]
                      ]
                    ]
                  ]
                ]
                [{member(enlace($Intermedio_{1}$,$C_{1}$),\lbrack enlace(winnipeg,20)\rbrack ),\\ruta\_ingenua($Intermedio_{1}$,calgary)}
                  [{\ligadura{$Intermedio_{1}$/winnipeg, $C_{1}$/20}\\ruta\_ingenua(winnipeg,calgary)}
                    [{conexion\_directa(winnipeg,calgary,$C_{2}$)}
                      [{nodo(winnipeg,$Enlaces_{2}$),\\member(enlace(calgary,$C_{2}$),$Enlaces_{2}$)}
                        [{\ligadura{$Enlaces_{2}$/\lbrack\rbrack}\\member(enlace(calgary,$C_{2}$),\lbrack\rbrack )}
                          [{false}, fallo
                          ]
                        ]
                      ]
                    ]
                    [{conexion\_directa(winnipeg,$Intermedio_{2}$,$C_{2}$),\\ruta\_ingenua($Intermedio_{2}$,calgary)}
                      [{nodo(winnipeg,$Enlaces_{2}$),\\member(enlace($Intermedio_{2}$,$C_{2}$),$Enlaces_{2}$),\\ruta\_ingenua($Intermedio_{2}$,calgary)}
                        [{\ligadura{$Enlaces_{2}$/\lbrack\rbrack}\\member(enlace($Intermedio_{2}$,$C_{2}$),\lbrack\rbrack ),\\ruta\_ingenua($Intermedio_{2}$,calgary)}
                          [{false}, fallo
                          ]
                        ]
                      ]
                    ]
                  ]
                  [{member(enlace($Intermedio_{1}$,$C_{1}$),\lbrack\rbrack ),\\ruta\_ingenua($Intermedio_{1}$,calgary)}
                    [{false}, fallo
                    ]
                  ]
                ]
              ]
            ]
          ]
        ]
        [{member(enlace($Intermedio_{0}$,$C_{0}$),\lbrack\rbrack ),\\ruta\_ingenua($Intermedio_{0}$,calgary)}
          [{false}, fallo
          ]
        ]
      ]
    ]
  ]
]
\end{forest}
}
    \caption{Árbol SLD de \texttt{ruta\_ingenua(edmonton, calgary)} (13 primeros niveles de un árbol infinito).}
\end{figure}
\end{landscape}

\newpage


%===============================================================================
% REPOSITORIO
%===============================================================================
\chapter{Repositorio}

El código fuente de esta actividad y el de las demás actividades del semestre está disponible en GitHub; el historial de versiones del repositorio sirve como soporte del trabajo realizado:

\begin{center}
    \url{https://github.com/1594-adrs/ProgIIIG1-Andres_David_Rincon_Salazar}
\end{center}


\end{document}
